% Abstract - written by M1 (Saksham Agrawal). 150-200 words; numbers come from the results.
\section*{Abstract}
\addcontentsline{toc}{section}{Abstract}
Autonomous drones fly most commercial missions as lists of waypoints, so they must hold altitude and track paths accurately despite wind and sensor noise. This project develops altitude and waypoint control for the DJI Mavic 2 Pro model in the Webots simulator, using a Python controller. We first measured the drone's control-sign conventions with a probe, then compared two cascaded controllers on the same attitude loop. Method A is a sample-style P/PD design with step waypoints. Method B adds a minimum-jerk trajectory, PID with anti-windup, a position--velocity cascade and feedforward; both methods share a Kalman velocity estimator. Six experiments covered altitude steps, square and 3D missions, crosswind, sensor noise over ten seeds and a roll-gain sweep, 42 runs in total. Method B reduced the steady altitude error from \R{E1}{baseline}{step2_ss_error_m}\,m to \R{E1}{improved}{step2_ss_error_m}\,m and the 3D path RMSE from \R{E3}{baseline}{path_rmse_m}\,m to \R{E3}{improved}{path_rmse_m}\,m, and it completed the crosswind mission, which Method A could not. The novelty is a measured finding: the sample's roll gain crashes the step-command controller, while smooth references make the system tolerant of that gain.

\textbf{Keywords:} quadrotor, PID control, trajectory generation, Kalman filter, Webots.
